% !TeX program = lualatex
% !TeX encoding = utf8
% !TeX spellcheck = uk_UA
% !BIB program = biber

\documentclass{PySCFBook}


\begin{document}

%% ========================================================
\chapter{Методи градієнтів та молекулярні властивості}\label{GradientMethods}
%% ========================================================

\section{Вступ: чому потрібні похідні енергії?}

Після розв’язання електронного рівняння Шредінґера та отримання електронної енергії $E$ відкривається можливість обчислення численних молекулярних властивостей. Найважливішою серед них є \emph{рівноважна молекулярна геометрія}. Обчислення структур молекул є цінним доповненням до експериментальних даних у галузях структурної хімії, таких як рентгенівська кристалографія, електронна дифракція та мікрохвильова спектроскопія.

Ключ до ефективного визначення рівноважних структур --- обчислення похідних потенціальної енергії за ядерними координатами. Ці похідні можна отримувати \emph{чисельно}, обчислюючи енергію в багатьох геометріях та визначаючи зміну енергії при варіюванні кожної ядерної координати. Однак \emph{методи градієнтів}, які визначають похідні енергії \emph{аналітично}, є значно швидшими та точнішими за чисельне диференціювання.

%% --------------------------------------------------------
\section{Похідні енергії та матриця Гессе}

З 1969 року, коли П. Пулай написав першу програму для аналітичного обчислення перших похідних \texttt{SCF}-енергій, методи градієнтів стали однією з найактивніше досліджуваних галузей сучасної квантової хімії. Спочатку застосовані до замкнутооболонкових \texttt{SCF}-розрахунків, вони згодом були узагальнені на відкриті оболонки (\texttt{RHF} та \texttt{UHF}). Аналітичні вирази також виведено для теорії функціоналу густини (\texttt{DFT}). Сьогодні аналітичні перші та другі похідні енергії доступні для більшості рівнів ab initio обчислень.

Для двоатомної молекули потенціальна енергія $E$ залежить лише від між’ядерної відстані $R$. Щоб знайти мінімум потенціалу (або будь-яку стаціонарну точку), потрібно знайти нуль похідної:
\begin{equation*}
\frac{dE}{dR} = 0.
\end{equation*}

Для багататомних молекул потенціальна енергія --- функція багатьох ядерних координат $q_i$. У рівноважній геометрії кожна сила $f_i$, що діє на ядро з боку електронів та інших ядер, має зникати:
\begin{equation*}
f_i = -\frac{\partial E}{\partial q_i} = 0.
\end{equation*}

Отже, рівноважну геометрію можна знайти, обчислюючи всі сили в заданій геометрії та перевіряючи, чи вони дорівнюють нулю. Якщо ні --- геометрію змінюють, доки не буде знайдено конфігурацію з нульовим градієнтом. Обчислювально сили не зникають ідентично, тому ітераційний пошук припиняють, коли величини сил менші за заданий поріг (наприклад, $10^{-5}$ Гартрі/Бор).

Зазвичай, щоб отримати кути зв’язку та торсійні кути з точністю $\pm 1^\circ$, а довжини зв’язку --- $\pm 0{,}002\,a_0$ (0,1 пм), усі сили мають бути меншими за $\sim 10$ пН. Такі критерії збіжності потребують від $N$ до $2N$ ітерацій, де $N$ --- число атомів у молекулі.

Перед оптимізацією геометрії необхідно обрати систему координат для представлення потенціальної поверхні. Вибір важливий, оскільки впливає на швидкість оптимізації. Часто використовують \emph{внутрішні координати} (довжини зв’язків, кути, торсійні кути).

Важливо пам’ятати, що експериментальні дані зазвичай стосуються \emph{усереднених} молекулярних величин. Тому при порівнянні з експериментом слід обчислювати \emph{вібраційно усереднені} довжини зв’язків та кути. Різниці між експериментальними та обчисленими рівноважними геометріями для серії споріднених сполук часто виявляються дуже стабільними. Наприклад, у \texttt{HF-SCF} дослідженні 30 органічних сполук з базисом 4-21G довжина зв’язку \ce{C-H} систематично була на $0{,}07\,a_0$ меншою за експериментальну. Така стабільність дозволяє створювати набори емпіричних поправок, що забезпечують абсолютну точність $\sim 0{,}02\,a_0$.

Нульовий градієнт характеризує \emph{стаціонарну точку} на поверхні, але не розрізняє мінімуми, максимуми та сідлові точки. Пошук дозволяє знаходити не лише рівноважну геометрію стабільної молекули, а й \emph{перехідний стан} хімічної реакції, що відповідає сідловій точці.

Для розрізнення типів стаціонарних точок необхідно розглядати \emph{другі похідні} енергії за ядерними координатами. Величини
\begin{equation*}
\frac{\partial^2 E}{\partial q_i \partial q_j}
\end{equation*}
утворюють \emph{матрицю Гессе}.

Мінімум (максимум) на одновимірній потенціальній кривій відповідає додатній (від’ємній) другій похідній. На багатовимірній поверхні мінімум (максимум) характеризується \emph{усіма додатними (від’ємними) власними значеннями матриці Гессе}. Перехідний стан (сідлова точка першого порядку) відповідає \emph{одному від’ємному} та всім іншим додатним власним значенням.

%% --------------------------------------------------------
\section{Аналітичні похідні та зв’язані збурені рівняння}

Існує багато алгоритмів пошуку стаціонарних точок на потенціальній поверхні. Їх можна класифікувати за використанням похідних:

\begin{itemize}
  \item \textbf{Тільки енергія} --- найповільніші, але корисні, якщо аналітичні похідні недоступні.
  \item \textbf{Енергія + перші похідні} --- значно ефективніші (майже на порядок). Швидкість збіжності зростає при наявності гарного початкового наближення до матриці Гессе (наприклад, з нижчого рівня ab initio).
  \item \textbf{Енергія + перші + другі похідні} --- найточніші та найшвидші.
\end{itemize}

У всіх випадках необхідно оптимізувати \emph{всі} ядерні координати, особливо для перехідних станів, де часткова оптимізація може дати некоректну сідлову точку.

Похідні енергії корисні не лише для геометрії. Елементи матриці Гессе --- це \emph{силові сталі} для частот нормальних мод у гармонічному наближенні (див. розділ 10.13). Треті, четверті та вищі похідні дають ангармонічні поправки до вібраційних частот (розділ 10.16).

Похідні не обмежуються ядерними координатами. Наприклад, похідні за компонентами електричного поля дають дипольні моменти. Змішані другі похідні (за ядерною координатою та компонентою поля) --- похідні дипольного моменту, що визначають інтенсивності в ІЧ-спектрах у гармонічному наближенні (розділ 10.10).

%% --------------------------------------------------------
\subsection{Обчислення аналітичних похідних}

Для обчислення аналітичних похідних енергії за ядерними координатами необхідно обчислювати похідні одно- та двоелектронних інтегралів над базисними функціями. Оскільки базисні функції центровані на ядрах, при диференціюванні інтегралів потрібні похідні самих базисних функцій за ядерними координатами.

Чи потрібні похідні коефіцієнтів розкладу (наприклад, молекулярних орбіталей), залежить від того, чи вони визначені варіаційно, та від порядку похідної енергії.

Розглянемо загальний коефіцієнт розкладу $c_j$. Перша похідна енергії:
\begin{equation*}
\frac{dE}{dq_i}
= \frac{\partial E}{\partial q_i}
+ \sum_j \frac{\partial E}{\partial c_j} \frac{\partial c_j}{\partial q_i}.
\end{equation*}

Для \emph{варіаційно визначених} коефіцієнтів член $\partial E / \partial c_j = 0$. Отже, для обчислення \emph{градієнта} (першої похідної) енергії \emph{не потрібні похідні коефіцієнтів}. У методах \texttt{HF} та \texttt{MCSCF} аналітичний градієнт потребує лише похідних одно- та двоелектронних інтегралів.

Друга похідна:
\begin{equation*}
\frac{d^2E}{dq_i^2}
= \frac{\partial^2 E}{\partial q_i^2}
+ \sum_j \left(
  \frac{\partial^2 E}{\partial q_i \partial c_j} \frac{\partial c_j}{\partial q_i}
  + \frac{\partial E}{\partial c_j} \frac{\partial^2 c_j}{\partial q_i^2}
\right).
\end{equation*}

Оскільки $\partial E / \partial c_j = 0$, друга похідна \emph{не потребує других похідних коефіцієнтів}, але потребує \emph{перших похідних} $\partial c_j / \partial q_i$. Вони обчислюються розв’язанням \emph{зв’язаних збурених рівнянь Хартрі-Фока} (\texttt{CPHF}) для методів з одним детермінантом (\texttt{SDCI}, \texttt{MP2}) або \emph{зв’язаних збурених рівнянь \texttt{MCSCF}} (\texttt{CPMCSCF}) для багатодетермінアントних методів (\texttt{MCSCF}, \texttt{MRCI}).

Загалом, для похідних енергії порядку $2n+1$ потрібні похідні варіаційних коефіцієнтів порядку $n$. Отже, третя похідна енергії потребує перших похідних коефіцієнтів.

Ефективний метод розв’язання \texttt{CPHF} розробив Дж. А. Попл у 1979 році, що зробило обчислення других похідних енергії практичними для \texttt{RHF} та \texttt{UHF}.

Для \emph{неваріаційних} коефіцієнтів (наприклад, у \texttt{MP2}) потрібні похідні коефіцієнтів. Їх отримують розв’язанням відповідних зв’язаних збурених рівнянь. Однак Н. К. Генді та Г. Ф. Шефер у 1984 році показали, що замість розв’язання всіх \texttt{CPHF}/\texttt{CPMCSCF} рівнянь достатньо розв’язати \emph{набагато меншу систему}. Це спрощення застосовне й до вищих похідних.

Для обчислення похідних інтегралів потрібні похідні базисних функцій за координатами їхніх центрів. Для гауссових орбіталей (\texttt{GTO}) похідні обчислюються аналітично та призводять до інших \texttt{GTO}. Наприклад, перша похідна s-\texttt{GTO} за $x_c$ дає p-\texttt{GTO}; похідна p-\texttt{GTO} $y_{100}$ --- s- та d-\texttt{GTO} (див. задачу 9.25).

%% --------------------------------------------------------
\section{Висновок}

Методи градієнтів --- основа сучасної обчислювальної хімії. Вони дозволяють:

\begin{itemize}
  \item Швидко та точно знаходити рівноважні геометрії та перехідні стани.
  \item Обчислювати частоти коливань, інтенсивності спектрів, поляризованість тощо.
  \item Уникати чисельного шуму та помилок скінченних різниць.
\end{itemize}

Аналітичні похідні, реалізовані через \texttt{CPHF} та \texttt{CPMCSCF}, разом із ефективними алгоритмами оптимізації, зробили можливими розрахунки великих молекул на високому рівні теорії.

\begin{table}[h!]
\centering
\caption{Зв’язок порядку похідних енергії з молекулярними властивостями (доповнення до табл.~\ref{tab:prop}).}\label{tab:grad}
\begin{tblr}{
  colspec={Q[c,m] Q[l,m]},
  cell{1}{1-2}={c=2}{c},
  hline{2} = {2pt}, hline{3} = {1pt}, hline{Z} = {1pt},
}
Порядок похідної & Властивість \\
1 & Сили на ядрах, градієнт \\
2 & Силові сталі, частоти коливань, гессіан \\
3 & Ангармонічні поправки \\
1 (за $R$) + 1 (за $E$) & Похідні дипольного моменту, ІЧ-інтенсивності \\
\end{tblr}
\end{table}

Таким чином, \emph{всі молекулярні властивості --- це похідні енергії}, а методи градієнтів --- універсальний інструмент їхнього аналітичного обчислення.

\end{document}
